% Λύση του 31ου προβλήματος της ενότητας «Άλυτα Προβλήματα».
\subsection*{Λύση Προβλήματος 31 --- Cobb--Douglas και οριακό προϊόν κεφαλαίου}

Για $L=100$ η παραγωγική συνάρτηση γίνεται
\[
  Q(K)=1{,}01\cdot100^{3/4}K^{1/4}.
\]
Θέτουμε
\[
  B=1{,}01\cdot100^{3/4}.
\]
Τότε $Q(K)=BK^{1/4}$.

\begin{alist}
  \item Εφόσον η τιμή πώλησης κάθε μονάδας παραγωγής είναι $1$, τα
        έσοδα είναι
        \[
          R(K)=Q(K)=BK^{1/4}.
        \]
        Το μεταβλητό κόστος χρήσης κεφαλαίου είναι
        \[
          C(K)=2K.
        \]
        Άρα η συνάρτηση κέρδους, πριν από το σταθερό κόστος εργασίας,
        είναι
        \[
          {\Pi(K)=BK^{1/4}-2K}.
        \]

  \item Το οριακό προϊόν του κεφαλαίου είναι
        \[
          Q'(K)=\frac{B}{4K^{3/4}}.
        \]
        Επειδή
        \[
          Q''(K)=-\frac{3B}{16K^{7/4}}<0
          \qquad (K>0),
        \]
        το οριακό προϊόν του κεφαλαίου μειώνεται καθώς αυξάνεται το
        κεφάλαιο. Αυτό είναι το φαινόμενο των φθινουσών οριακών
        αποδόσεων.

  \item Παραγωγίζουμε τη συνάρτηση κέρδους:
        \[
          \Pi'(K)=\frac{B}{4K^{3/4}}-2.
        \]
        Για μέγιστο ζητούμε $\Pi'(K)=0$, δηλαδή
        \[
          \frac{B}{4K^{3/4}}=2.
        \]
        Άρα
        \begin{align*}
          K^{3/4}&=\frac{B}{8},\\
          K^*&=\left(\frac{B}{8}\right)^{4/3}.
        \end{align*}
        Επειδή
        \[
          \Pi''(K)=-\frac{3B}{16K^{7/4}}<0,
        \]
        το σημείο αυτό δίνει μέγιστο κέρδος.

        Επομένως,
        \[
          {
            K^*=
            \left(
              \frac{1{,}01\cdot100^{3/4}}{8}
            \right)^{4/3}
            \approx6{,}33
          }.
        \]
        Η συνθήκη
        \[
          \frac{B}{4K^{3/4}}=2
        \]
        λέει ότι το οριακό έσοδο που προσφέρει μία επιπλέον μονάδα
        κεφαλαίου είναι ίσο με το οριακό κόστος χρήσης της.

  \item Για $K^*\approx6{,}33$ έχουμε
        \[
          Q(K^*)=B(K^*)^{1/4}\approx50{,}67.
        \]
        Επίσης,
        \[
          C(K^*)=2K^*\approx12{,}67.
        \]
        Άρα
        \[
          \Pi(K^*)=Q(K^*)-2K^*
          \approx50{,}67-12{,}67
          \approx38{,}00.
        \]
        Συνεπώς,
        \[
          {Q(K^*)\approx50{,}67},
          \qquad
          {\Pi(K^*)\approx38{,}00}.
        \]
\end{alist}
